\section*{Tóm tắt}
\thispagestyle{empty}

Đề tài “Phát triển hệ thống quản lý thư viện trực tuyến” được thực hiện nhằm đáp ứng nhu cầu chuyển đổi số trong lĩnh vực quản lý và khai thác tài nguyên thư viện. Trong bối cảnh các thư viện truyền thống đang dần chuyển mình sang môi trường số, việc xây dựng một hệ thống quản lý hiện đại, hiệu quả và có khả năng hỗ trợ người dùng tốt hơn là một yêu cầu cấp thiết. Đề tài tập trung nghiên cứu, thiết kế và triển khai một hệ thống thư viện dưới dạng ứng dụng Web (Web Application), hướng đến các thư viện có quy mô vừa và nhỏ như thư viện trường học, thư viện khoa hoặc trung tâm học liệu. \\

Hệ thống được xây dựng nhằm quản lý đa dạng các loại tài nguyên thư viện bao gồm sách in, tạp chí và các tài nguyên số như E-book, tài liệu PDF. Đồng thời, hệ thống không chỉ dừng lại ở việc hỗ trợ các chức năng quản lý cơ bản mà còn hướng đến nâng cao trải nghiệm người dùng thông qua việc tích hợp các công nghệ hiện đại, đặc biệt là trí tuệ nhân tạo (AI). \\

Trong hệ thống, các đối tượng người dùng chính bao gồm bạn đọc (reader), thủ thư (librarian) và quản trị viên (admin). Mỗi vai trò sẽ có những chức năng và nhu cầu sử dụng khác nhau, được hệ thống hỗ trợ một cách phù hợp.\\

Đối với vai trò bạn đọc (reader), hệ thống cung cấp các chức năng hỗ trợ tìm kiếm và tra cứu tài liệu một cách linh hoạt thông qua từ khóa hoặc các công cụ tìm kiếm thông minh. Bên cạnh đó, hệ thống còn tích hợp cơ chế gợi ý tài liệu dựa trên lịch sử đọc và hành vi người dùng, giúp người đọc nhanh chóng tiếp cận được những tài liệu phù hợp với nhu cầu học tập và nghiên cứu. Người dùng có thể quản lý lịch sử mượn sách, theo dõi trạng thái mượn – trả và tương tác với hệ thống thông qua chatbot để được hỗ trợ giải đáp các thắc mắc một cách nhanh chóng.\\

Đối với vai trò thủ thư (librarian), hệ thống hỗ trợ quản lý danh mục tài nguyên, cập nhật thông tin tài liệu và xử lý các hoạt động mượn – trả. Ngoài ra, thủ thư có thể theo dõi tình trạng tài nguyên, kiểm soát dữ liệu và đảm bảo quá trình vận hành thư viện diễn ra một cách hiệu quả. Việc ứng dụng hệ thống giúp giảm thiểu các thao tác thủ công, từ đó nâng cao năng suất làm việc và hạn chế sai sót trong quá trình quản lý.\\

Đối với quản trị viên (admin), hệ thống cung cấp quyền quản lý toàn bộ hoạt động, bao gồm quản lý tài khoản người dùng, cấu hình hệ thống và giám sát dữ liệu. Bên cạnh đó, admin có thể theo dõi các thống kê liên quan đến hoạt động mượn – trả, từ đó hỗ trợ việc đưa ra các quyết định quản lý và phát triển nguồn tài nguyên thư viện một cách hợp lý.\\

Một điểm nổi bật của hệ thống là việc tích hợp các chức năng trí tuệ nhân tạo ở mức cơ bản như gợi ý tài liệu, tìm kiếm bằng ngôn ngữ tự nhiên và chatbot hỗ trợ người dùng. Các chức năng này giúp nâng cao khả năng tương tác giữa người dùng và hệ thống, đồng thời cải thiện đáng kể trải nghiệm sử dụng so với các hệ thống thư viện truyền thống.\\

Về mặt kiến trúc, hệ thống được xây dựng theo mô hình ứng dụng web với các thành phần chính bao gồm giao diện người dùng, hệ thống xử lý nghiệp vụ và cơ sở dữ liệu. Cách tổ chức này giúp đảm bảo tính linh hoạt, dễ mở rộng và thuận tiện trong việc bảo trì, nâng cấp hệ thống trong tương lai.\\

Thông qua việc xây dựng hệ thống quản lý thư viện trực tuyến, đề tài hướng đến mục tiêu tối ưu hóa quy trình quản lý tài nguyên, nâng cao hiệu quả vận hành thư viện và mang lại trải nghiệm tốt hơn cho người dùng trong việc tìm kiếm và khai thác tài liệu. Đồng thời, đề tài cũng góp phần minh họa khả năng ứng dụng của các công nghệ hiện đại, đặc biệt là trí tuệ nhân tạo, trong việc phát triển các hệ thống thông tin phục vụ lĩnh vực giáo dục và nghiên cứu.\\

\clearpage
\pagenumbering{arabic}
